\documentclass[11pt,letterpaper]{article}
\usepackage[margin=1in,headheight=14pt]{geometry}
\usepackage[T1]{fontenc}
\usepackage{lmodern}
\usepackage{amsmath,amssymb,amsthm,mathtools}
\usepackage{microtype,booktabs,xcolor,fancyhdr}
\definecolor{ink}{HTML}{19354A}
\usepackage[colorlinks=true,linkcolor=ink,urlcolor=ink,citecolor=ink]{hyperref}
\hypersetup{pdftitle={Expansion and random walks: exact examples},pdfauthor={Akhil Ravuri}}
\newtheorem{proposition}{Proposition}
\newcommand{\TV}{\operatorname{TV}}
\pagestyle{fancy}\fancyhf{}
\fancyhead[L]{\small\textsc{Expansion and random walks}}
\fancyhead[R]{\small Akhil Ravuri}\fancyfoot[C]{\thepage}
\renewcommand{\headrulewidth}{0.3pt}
\setlength{\parskip}{0.25em}\setlength{\emergencystretch}{1em}
\title{\vspace{-1.5em}\LARGE\bfseries Expansion and random walks:\\ exact examples}
\author{Akhil Ravuri}\date{}
\begin{document}
\maketitle\vspace{-1.5em}
\begin{abstract}
We give elementary proofs of edge-expansion formulas for complete graphs and
cycles, and exact total-variation formulas for ordinary and lazy random walks
on complete bipartite graphs. Disjoint cliques show why laziness cannot repair
disconnection. Rational computations in the accompanying repository provide
finite implementation checks. This is an expository note; no novelty is claimed.
\end{abstract}

\section{Cuts and stationary distributions}
Let $G=(V,E)$ be a finite simple graph on $n\ge2$ vertices. Write
$\partial S=E(S,V\setminus S)$ and use the unnormalized edge expansion
\[
 h(G)=\min_{0<|S|\le n/2}\frac{|\partial S|}{|S|}.
\]
This is the convention of Hoory, Linial, and Wigderson~\cite[Definition 2.1]{HLW}.
It differs from degree-normalized conductance.

\begin{proposition}\label{prop:cuts}
For $n\ge2$, $h(K_n)=\lceil n/2\rceil$. For $n\ge3$,
$h(C_n)=2/\lfloor n/2\rfloor$.
\end{proposition}
\begin{proof}
A set of $s$ vertices in $K_n$ has $s(n-s)$ boundary edges. Its ratio is
$n-s$, minimized at $s=\lfloor n/2\rfloor$. In a cycle every nonempty proper
vertex set has at least two boundary edges, so every eligible set has ratio
at least $2/\lfloor n/2\rfloor$. A consecutive block of
$\lfloor n/2\rfloor$ vertices attains this bound.
\end{proof}

Now assume $G$ is $d$-regular with $d\ge1$. The ordinary transition matrix is
$P=A/d$, where $A$ is the adjacency matrix. The lazy transition is
$L=(I+P)/2$. Distributions are row vectors, so $p_{t+1}=p_tP$ or $p_tL$.
The uniform distribution $\pi(v)=1/n$ is stationary for both walks: every
column of each transition matrix sums to one. This remains true when the
graph is disconnected.

For distributions $p,q$ on $V$, define
\[
 \TV(p,q)=\frac12\sum_{v\in V}|p(v)-q(v)|.
\]
Stationarity says that a walk \emph{started} uniformly remains uniform.
It does not say that a walk started at one vertex converges to uniform.
The next examples separate periodicity from disconnection.

\newpage
\section{A bipartite example with an exact mixing formula}
Let $G=K_{d,d}$, $d\ge1$, with parts $U,W$, and start at a vertex $x\in U$.
Let $u_U,u_W$ denote the uniform distributions supported on the respective
parts, extended by zero elsewhere. Put
$\pi=(u_U+u_W)/2$ and $\delta_x(x)=1$.

\begin{proposition}\label{prop:bipartite}
For the ordinary walk on $K_{d,d}$,
\[
 \TV(\delta_xP^t,\pi)=\frac12 \qquad(t\ge1).
\]
For the lazy walk, $\TV(\delta_x,\pi)=1-1/(2d)$, and
\[
 \TV(\delta_xL^t,\pi)=\frac{d-1}{d}\,2^{-t}
 \qquad(t\ge1).
\]
\end{proposition}
\begin{proof}
After one ordinary step the distribution is $u_W$. Later steps alternate
between $u_U$ and $u_W$. Either distribution differs from $\pi$ by
$1/(2d)$ at each of $2d$ vertices, giving distance $1/2$.

For the lazy walk, decompose the initial distribution as
\[
 \delta_x=\pi+\frac12(u_U-u_W)+(\delta_x-u_U).
\]
Multiplication by $P$ swaps $u_U$ and $u_W$, so
$(u_U-u_W)L=0$. Also $(\delta_x-u_U)P=0$: both $\delta_xP$
and $u_UP$ equal $u_W$. Hence
\[
 \delta_xL^t=\pi+2^{-t}(\delta_x-u_U),\qquad t\ge1.
\]
The sum of absolute values of $\delta_x-u_U$ is $2(d-1)/d$.
Taking half the absolute sum proves the formula. At time zero, the distance
from a point mass to uniform on $2d$ vertices is $1-1/(2d)$.
\end{proof}

The restriction $t\ge1$ matters: the imbalance between the two parts
disappears in the first lazy step, whereas the remaining imbalance within
$U$ decays by a factor of two per step. For $d=1$ the latter imbalance is
zero, so the lazy walk is uniform after one step.

For $K_{3,3}$, the time-eight lazy distance is
$(2/3)2^{-8}=1/384$, while the ordinary distance stays $1/2$.
More generally, for $d\ge2$ and a tolerance $0<\varepsilon<1/2$, the first
integer time at which the lazy distance is at most $\varepsilon$ is
\[
 t_\varepsilon=\left\lceil\log_2\frac{d-1}{d\varepsilon}\right\rceil.
\]
Indeed, the time-zero distance exceeds $1/2$, and for $t\ge1$ the desired
inequality is equivalent to $2^t\ge(d-1)/(d\varepsilon)$.
On $K_{3,3}$, tolerance $1/100$ is first reached at time seven:
the distances at times six and seven are $1/96$ and $1/192$, respectively.

On a general regular bipartite graph the ordinary walk still remains on
one part at each time, giving distance \emph{at least} $1/2$ from uniform;
the exact equality above uses the complete bipartite structure.

These examples show why a positive edge-expansion value alone cannot force
an ordinary walk to converge. For example $K_{2,2}=C_4$ has expansion one,
yet its ordinary walk alternates between the two parts forever.

\newpage
\section{What laziness cannot fix}
Let $G$ be the disjoint union of $r\ge2$ copies of $K_d$, with $d\ge2$.
Start at $x$ in one component $H$. Here the graph degree is $d-1$.
Write $u_H$ for uniform measure on that component and let $\pi$ be uniform
on all $rd$ vertices.

\begin{proposition}\label{prop:disconnected}
The graph $G$ has $h(G)=0$. Any walk that stays in its starting component
has distance at least $1-1/r$ from $\pi$. For the lazy random walk,
\[
 \delta_xL^t=u_H+
 \left(\frac{d-2}{2(d-1)}\right)^t(\delta_x-u_H),\qquad t\ge1,
\]
and therefore $\lim_{t\to\infty}\TV(\delta_xL^t,\pi)=1-1/r$.
\end{proposition}
\begin{proof}
The set $V(H)$ has size $d\le rd/2$ and no boundary edges, proving
$h(G)=0$. If a distribution $p$ is supported on $H$, the absolute
differences outside $H$ sum to $1-1/r$. Inside $H$, their sum is at least
$|p(H)-\pi(H)|=1-1/r$. Thus $\TV(p,\pi)\ge1-1/r$.

On $H$, $P=(J-I)/(d-1)$, where $J$ is the all-ones matrix.
Since $(\delta_x-u_H)J=0$, multiplication by $L$ scales
$\delta_x-u_H$ by $(d-2)/(2(d-1))$, while fixing $u_H$.
Iteration proves the displayed formula. This factor belongs to $[0,1)$,
so the distribution converges to $u_H$. Finally,
$\TV(u_H,\pi)=1-1/r$ by direct summation.
\end{proof}

\subsection*{Exact computations and scope}
The existing experiment enumerates all eligible cuts and evolves walk
probabilities using rational arithmetic. Walks start at vertex 0.
The following entries are the repository's recorded time-eight values.
\begin{center}
\begin{tabular}{lrrr}
\toprule
Graph & $h(G)$ & Lazy TV & Ordinary TV\\
\midrule
$K_4$ & $2$ & $1/8748$ & $1/8748$\\
$C_6$ & $2/3$ & $2187/32768$ & $1/2$\\
$K_{3,3}$ & $5/3$ & $1/384$ & $1/2$\\
$2K_3$ & $0$ & $1/2$ & $1/2$\\
\bottomrule
\end{tabular}
\end{center}
Reproduce from the repository root with
\texttt{python -m experiments.expansion\_small\_cases}.
Formula checks are in \texttt{tests/test\_walk\_formulas.py}.
The propositions have proofs independent of these finite checks.
No infinite expander family, spectral inequality, or high-dimensional
expansion result is established here. For broader context see~\cite[Sections 2--3]{HLW}.

\begin{thebibliography}{1}
\bibitem{HLW} S. Hoory, N. Linial, and A. Wigderson,
\emph{Expander graphs and their applications}, Bull. Amer. Math. Soc.
\textbf{43} (2006), 439--561.
\href{https://doi.org/10.1090/S0273-0979-06-01126-8}{doi:10.1090/S0273-0979-06-01126-8}.
\end{thebibliography}
\end{document}
